\documentclass[11pt]{article}
\usepackage[T1]{fontenc}
\usepackage{textcomp}
\usepackage[margin=0.75in]{geometry}
\usepackage{amsmath,amssymb,amsthm}
\usepackage{array,booktabs}
\usepackage{hyperref}
\setlength{\parindent}{0pt}
\newtheorem{theorem}{Theorem}
\theoremstyle{definition}
\newtheorem{definition}{Definition}
\title{Flow Proof Artifact: GroupHom-inverse-product}
\author{Generated by \texttt{flow doc proof}}
\date{Flow Proof Book}
\begin{document}
\maketitle
Homomorphisms send inverse products to one.\\[0.5em]
\textbf{Source.} dummit-foote — *Abstract Algebra*, §1.6\\[0.5em]
\bigskip
\noindent{\large \textbf{Derived fact 1.} \textit{f(g) * f(inv(g)) = 1} \hfill GroupHom \cdot multiplication \cdot products with inverses map to one}\par
\smallskip\noindent\textit{Coordinate: GroupHom \cdot multiplication \cdot products with inverses map to one}\par
\smallskip\noindent\textit{Source: dummit-foote}\par
\smallskip\noindent\textit{Built on: products map to products, for multiplication on GroupHom, inverses map to inverses, for inverse on GroupHom, right inverse recovers the identity, for inverse on Group}\par
\medskip
\noindent\textbf{Goal.} f(g) * f(inv(g)) = 1\par
\noindent$\displaystyle \forall f \in GroupHom \forall g \in Group\quad f(g) * f(inv(g)) = 1$\par
\medskip
\noindent
\begin{tabular}{@{} >{\raggedright\arraybackslash}p{0.06\textwidth} >{\raggedright\arraybackslash}p{0.47\textwidth} >{\raggedleft\arraybackslash}p{0.41\textwidth} @{}}
\textbf{\#} & \textbf{Proof} & \textbf{Mathematics} \\
\midrule
\textbf{1.} & We prove that products with inverses map to one for multiplication on GroupHom. & \\[0.45em]
\textbf{2.} & We invoke the definitional clause governing multiplication on GroupHom: products map to products, for multiplication on GroupHom (instantiated for f, g, inv(g)). & \\[0.45em]
\textbf{3.} & We invoke the derived fact governing inverse on GroupHom: inverses map to inverses, for inverse on GroupHom (instantiated for f, g). & \\[0.45em]
\textbf{4.} & We invoke the derived fact governing inverse on Group: right inverse recovers the identity, for inverse on Group (instantiated for g). & \\[0.45em]
\textbf{5.} & From step 2, step 3, and step 4, this implies f(g) times f(inv(g)) equals 1. Hence proven. & $\displaystyle f(g) * f(inv(g)) = 1$ \\[0.45em]
\bottomrule
\end{tabular}
\label{thm:GroupHom-multiplication-products_with_inverses_map_to_one}
\par\medskip\hrule
\end{document}
